\documentclass{article}

\usepackage[margin=1in]{geometry}
\usepackage{amsmath}
\usepackage{graphicx}
\usepackage{booktabs}
\usepackage{hyperref}

\title{Analysis of Parallelization Techniques on k-mer Counting for Genomic Sequences}

\author{Alex McGinn}
\date{\today}

\begin{document}

\maketitle

%========================================================== Introduction
\section{Introduction}

\vspace{\baselineskip}

In genomic sequence analysis, genetic information is represented by long sequences containing millions or even billions of nucleotide 
bases, primarily adenine (A), cytosine (C), guanine (G), and thymine (T). To compare genomic samples, scientists need methods that can 
identify meaningful similarities even when two sequences are not exact matches. This is important because even closely related organisms 
or samples can contain mutations, insertions, deletions, and other sequence variations. Although genomic sequences may differ, 
biologically meaningful similarities can often be detected through shared local sequence patterns such as k-mers \cite{moeckel2024kmer}.

\vspace{\baselineskip}

One method for identifying these local patterns is k-mer analysis. A sliding window of length \(k\) is moved across a genomic sequence, 
producing every contiguous substring of length \(k\). These substrings are called k-mers, and the number of times each unique k-mer occurs 
can then be counted. The resulting k-mer frequencies provide a compact way to characterize and compare genomic sequences without requiring 
the sequences to match in their entirety. Larger values of \(k\) provide more sequence-specific information, while smaller values provide 
greater sensitivity to local similarities \cite{lees2016seer}. For example, if two genomic regions contain the same 15-mer, this represents 
a more specific local match than sharing a much shorter k-mer. However, a single base difference can disrupt several overlapping large k-mers, 
while shorter k-mers are less sensitive to an individual mismatch.

\vspace{\baselineskip}

When analyzing a genetic sample, we do not obtain the entire genome in one clean pass. Instead, sequencing produces many shorter reads from the sample,
which often overlap and contain repeated coverage of the same genomic regions. As a result, the same k-mer may appear several times across the reads.
Because the complete genome may not already be known, we cannot always determine its total size directly from the sequencing output. However, as Marcais
and Kingsford explain, k-mer frequency statistics can be used to estimate sequencing coverage and subsequently ``derive an estimate of the genome size
from (c) and the total length of the reads'' \cite{marcais2011jellyfish}.

\vspace{\baselineskip}

Although, as k increases, the number of possible k-mers grows exponentially, leading to an increased amount of memory required to 
store and process k-mer counts. Alternatively, if k is shorter, we may have many matches, which would require the counter variable to 
have a larger range. 
For large genomic datasets, these computations become expensive, but there are various methods to accelerate this 
process using parallelization. The aim of this project was therefore to analyze the effectiveness of CPU/GPU parallelism for these  
types of tasks.

%========================================================== Background
\section{Background}

\vspace{\baselineskip}

To process a set of k-mer (determined by the k value; e.g. k = 4) we first move a sliding window of size k across the entire sequence and extract every k-mer substring.  
The first benefit is that we can now compare genomic sequences without simply comparing their complete ordering. Another benefit is that when we 
break a DNA sequence into many smaller subsequences of a known size, the condensed format allows us to perform computations in parallel, 
leading to a needed speedup. Although, as the input dataset size and k-value increase, the ability to store the data necessary in main memory becomes 
increasingly difficult.  

\vspace{\baselineskip}

To address this problem, one strategy for storing k-mers is to index an array or a similar data structure for each possible sequence. 
This allows us to predict the amount of main memory necessary to store all possible combinations for a dataset and set k-mer value.  
For example, a dataset could have a dedicated array that contains all possible k-mer permutation for k-value set equal to five.  
We know whether this would be possible or not because the number of possible permutations of such a string with 4  
possible bases is:
 
\[
4^k
\]

e.g. in a 5 k-mer example:

\[
4^5 = 1024 \text{ possible subsequences}
\]

\vspace{\baselineskip}

We could then also have an array that holds the count of each k-mer in the dataset. Therefore, we know that we  
need 1024 representations of the four possible bases. We can then also predict the size of the array used to count those  
occurrences. To determine the size necessary, we then need to consider the maximum possible number of occurences  
a k-mer can have. This will be dictated by frequency, which we cannot outright predict. However, we  
can be certain of the total number of possible occurrences based on the dataset size. That is because the maximum  
value will always be equal to how many possible substrings there are, determined by:

\vspace{\baselineskip}

\[
n - k + 1
\]

\vspace{\baselineskip}

e.g. in a DNA Sequence of length 66,000 with k-value equal to 5:

\vspace{\baselineskip}

\[
66000 - 5 + 1 = 65996
\]

\vspace{\baselineskip}

Which means that we need array elements that go up to at least 65996. In this instance, we would know we now need  
elements larger than an unsigned short, which would reach capacity at 65,535

\vspace{\baselineskip}

\subsection{Similar Work}

\vspace{\baselineskip}

One similar project proposed by Marcais and Kingsford \cite{marcais2011jellyfish} introduced Jellyfish, a parallel k-mer counting 
algorithm designed to improve performance and memory efficiency when processing large genomic datasets. Jellyfish utilizes a
multithreaded, lock-free hash table that allows concurrent updates through atomic operations rather than conventional
locking mechanisms. The implementation also leverages compact data representations to reduce memory requirements. Their work 
demonstrates the importance of synchronization overhead and memory-efficient data structures in parallel k-mer
counting, both of which will be considered when evaluating the parallelization strategies proposed in this project. 

\vspace{\baselineskip}

Deorowicz et al. \cite{deorowicz2015kmc2} introduced KMC 2, a parallel k-mer counting algorithm that focuses on reducing memory 
consumption and disk I/O overhead when processing large genomic datasets. Unlike approaches that rely primarily on concurrent hash-table 
updates, KMC 2 partitions sequence data into smaller groups using signatures derived from minimizers. These partitions can then be processed 
using parallel sorting and counting operations, reducing the amount of data that must be held in memory simultaneously. The authors also 
introduced $(k,x)$-mers to reduce redundant processing of overlapping k-mers. This approach highlights the importance of memory management, 
workload distribution, and I/O overhead, which are relevant considerations for the parallelization methodology proposed in this project.


%========================================================== Sequential Implementation

\section{Sequential Implementation}

\vspace{\baselineskip}

\subsection{Dataset and FASTA Input Processing}

\vspace{\baselineskip}

The initial sequential implementation was developed in C++ and tested using the Escherichia coli 
K-12 MG1655 reference genome, obtained from the National Center for Biotechnology Information (NCBI).
The reference genome contains 4,641,652 nucleotide bases and is provided in FASTA format \cite{ncbi_ecoli}.

\vspace{\baselineskip}

FASTA files contain base pairs separated by header lines beginning with the `\texttt{>}' character.
The implementation reads the file line by line, concatenating base pairs belonging to the same sequence
and storing each completed sequence in a string.

\vspace{\baselineskip}

The runtime is evaluated both before the FASTA loading operation is performed and before the sequential counting benchmark begins, allowing the
execution time of the counting algorithm to be measured independently of file input operations. 

\vspace{\baselineskip}

\subsection{Base Encoding}

\vspace{\baselineskip}

Each base is represented using a two-bit integer encoding, where A, C, G, and T correspond to the values
0, 1, 2, and 3, respectively. Since each base requires two bits, a k-mer of length $k$ can be represented
using $2k$ bits. The implementation then stores the encoded k-mer in an unsigned 64-bit integer using bit manipulation, 
which allows a maximum k-mer length of $k=32$.

\vspace{\baselineskip}

A count \texttt{vector} is allocated with $4^k$ elements, corresponding to every possible k-mer of length $k$.
Each element is a \texttt{uint64\_t}, resulting in a total allocation of $8 \times 4^k$ bytes. As each
k-mer is processed, its two-bit representation is shifted into the encoded integer. The resulting integer
value is then used directly as the index into the count \texttt{vector}, allowing the corresponding k-mer count
to be incremented without requiring a hash-table lookup.

\vspace{\baselineskip}

Many genetic sequences will have bases that are not read correctly, which is represented in the file as an \texttt{N}. 
Entries that contain this character representing an invalid reading will not be used in the final count.

\subsection{Sliding Window and Frequency Counting}

The sequential algorithm processes each sequence using a sliding window of length $k$. Rather than
constructing a new substring at every position, the implementation maintains an integer containing the
encoded bases of the current window.

For each new base, the encoded integer is shifted left by two bits, and the encoded value of the new
nucleotide is inserted using a bitwise OR operation. A bitmask is then applied to remove any bits extending
beyond the current k-mer, ensuring that only the most recent $k$ bases are retained.

Once $k$ valid nucleotides have been processed, the resulting integer is used as an index into a frequency
counting array. Each array entry stores the number of occurrences of its corresponding k-mer. This allows
the frequency counts to be updated directly without requiring a hash table or additional lookup operations.

\newpage

\subsection{Time \& Memory Complexity}

\vspace{\baselineskip}

The main components of the applications are reading the sequences from the file, initializing the counter array, and processing and incrementing each k-mer. 
Initializing the count of the array to default values takes $O(4^k)$ time, while processing and incrementing each k-mer only takes $O(n)$ time. Since reading 
from a file also takes linear time with respect to the number of bases, the overall time complexity is $O(4^k + n)$.

\vspace{\baselineskip}

%========================================================== Parallel Methodology
\section{Parallelization Methodology}

\vspace{\baselineskip}

The objective of this project is to analyze the effects of parallelization on k-mer counting for genomic sequences. 
The algorithm will then be written using OpenMP for parallelization on multiple CPU cores. Another implementation will utilize CUDA  
for GPU parallelization. MPI may also be evaluated if suitable distributed computing resources are available.  
There is a possibility of using an old desktop and Amazon EC2 instances to provide this additional resource. 

\vspace{\baselineskip}

Experiments will vary the dataset size and k-mer size. The performance will then be evaluated in standard methods. Most  
notably in speedup compared to the serial implementation, memory consumption, and the behavior of the implementation when scaling  
the dataset and k-mer size. 






%========================================================== Document End

\bibliographystyle{plain}
\bibliography{references}

\end{document}
%====================================================================================================================== EOF  


